\subsection{测度的例子}

例 1. --- 局部紧空间上的测度。

下面的命题表明，本章的理论包含 Ch. IV 的理论。在陈述中，“测度”一词及记号 \(\mathcal{M}(\mathrm{T};\mathbf{C})\) 都取前面各章中的意义。

命题 3. --- 设 T 是局部紧空间，\(\mu\) 是 T 上的一个测度。用 W( \(\mu\) ) 表示使 T 的每个紧子集 K 对应诱导测度 \(\mu_{K}\) 的映射。则 W( \(\mu\) ) 是 T 上的一个前测度，有 W(\textbar{} \(\mu\) \textbar) = \textbar W( \(\mu\) )\textbar，且线性映射 W: \(\mu \mapsto\) W( \(\mu\) ) 是空间 \(\mathcal{M}(T;C)\) 到 T 上前测度所成空间 \(\mathcal{P}(T;C)\) 上的一个双射。此外，若 \(\mu\) 是正的，则 \(\mu^{\bullet} = (W(\mu))^{\bullet}\) 。

显然 \(\mathrm{W}(\mu)\) 是一个前测度 (Ch. V, §7, No.~2, Prop. 4)，且映射 W 是线性的。关系 \(\mathrm{W}(\mu)=0\) 表示 \(\mu\) 在 T 的每个紧集上诱导测度 0；于是 \(\mu(f)=0\) 对 \(f\in\mathcal{K}(\mathrm{T};\mathbf{C})\) 成立，从而 \(\mu=0\) ，这证明 W 是单射。剩下证 W 是满射。由于每个前测度都是正前测度的线性组合，只需对每个正前测度 w 构造一个正测度 \(\mu\) 使 \(w=\mathrm{W}(\mu)\) 。设 \(f\in\mathcal{K}(\mathrm{T})\) 是给定的函数，L 是包含 f 的支集的一个紧集；由诱导测度的定义，数 \(w_{\mathrm{L}}(f_{\mathrm{L}})\) 与 L 的选取无关，故可令 \(\mu(f)=w_{\mathrm{L}}(f_{\mathrm{L}})\) ；于是 \(\mu\) 是 \(\mathcal{K}(\mathrm{T})\) 上的一个正线性形式，即一个正测度。我们来验证 \(w = \mathrm{W}(\mu)\) ；首先，关系式 \(\mu^{\bullet}(f) = w_{\mathrm{L}}^{\bullet}(f_{\mathrm{L}})\) 可推广到 \(f\) 是有限的正上半连续函数且在 \(\mathbf{L}\) 之外为零的情形。因为，设 \(\mathbf{M}\) 是 \(\mathbf{L}\) 的一个紧邻域，\(\mathcal{H}\) 是 \(\mathbf{T}\) 上支集含于 \(\mathbf{M}\) 且 \(\geqslant f\) 的连续函数所成的（递减定向）集合。则 (Ch. IV, §4, No.~4, Prop. 5 的 Cor. 2)

\[
\mu^ {\bullet} (f) = \inf _ {h \in \mathscr {H}} \mu (h) = \inf _ {h \in \mathscr {H}} w _ {\mathrm{M}} (h _ {\mathrm{M}}) = w _ {\mathrm{M}} ^ {\bullet} (f _ {\mathrm{M}}),
\]

另一方面，\(w_{\mathrm{M}}^{\bullet}(f_{\mathrm{M}}) = w_{\mathrm{L}}^{\bullet}(f_{\mathrm{L}})\) ，因为 \(f_{M}\) 在 M-L 上为零 (Ch. V, §7, No.~1, Prop. 1)。特别地，若把 f 取为 \(\mathcal{K}_{+}(\mathrm{L})\) 中元素经 0 延拓所得的函数，则由诱导测度的定义，此式表明 \(\mu_{L} = w_{L}\) ，从而确实有 \(\mathrm{W}(\mu) = w\) 。

若 \(\mu\) 是正的，则

\[
\mu^ {\bullet} (f) = \sup _ {K} \mu^ {\bullet} (f \varphi_ {K}) = \sup _ {K} \mu_ {K} ^ {\bullet} (f _ {K}) = (W (\mu)) ^ {\bullet} (f)
\]

对每个 \(f \in \mathcal{F}_{+}(\mathrm{T})\) 成立 (Ch. V, §1, Def. 1 与 §7, Prop. 1)。关系 \(|\mathrm{W}(\mu)| = \mathrm{W}(|\mu|)\) 是明显的 (Ch. IV, §5, No.~7, Lemma 3)。

证毕。

当 T 局部紧时，今后我们将借助双射 W 把空间 \(\mathcal{M}(\mathrm{T};\mathbf{C})\) 与 \(\mathcal{P}(\mathrm{T};\mathbf{C})\) 等同起来。

例 2. --- 拓扑空间上具有紧支集的测度。

引理 2. --- 设 T 是拓扑空间，L 是 T 的紧子集，\(\lambda\) 是 L 上的一个正测度。则存在 T 上唯一的正测度 \(\mu\) ，使得对每个函数 \(f \in \mathcal{F}_{+}(\mathrm{T})\) 有

\[
\mu^ {\bullet} (f) = \lambda^ {\bullet} (f _ {L}).\tag{2}
\]

令 \(p(f) = \lambda^{\bullet}(f_{\mathrm{L}})\) （对每个 \(f \in \mathcal{F}_{+}(\mathrm{T})\)），我们来证 Prop. 2 b) 的条件 1) 与 2) 得到满足。第二个条件显然满足：因为 \(p(f) = p(f\varphi_{\mathrm{K}})\) 只要 \(\mathrm{K}\) 包含 \(\mathrm{L}\) 。若 \(\mathrm{K} \subset \mathrm{T}\) 紧，且 \(h \in \mathcal{F}_{+}(\mathrm{K})\) ，则

\[
p _ {\mathrm{K}} (h) = p \left(h ^ {0}\right) = \lambda^ {\bullet} \left(h ^ {0} \mid \mathrm{L}\right).
\]

但 \(h^0 | \mathbf{L}\) 是 \(h_{\mathbf{K} \cap \mathbf{L}}\) 到 \(\mathbf{L}\) 上的 0 延拓：故最后的表达式等于 \((\mu_{\mathbf{K}})^{\bullet}(h)\) ，其中 \(\mu_{\mathbf{K}}\) 是 \(\lambda | \mathbf{K} \cap \mathbf{L}\) 在 \(\mathbf{K} \cap \mathbf{L}\) 到 \(\mathbf{K}\) 的单射下的像 (Ch. V, §6, No.~2, Prop. 2 与 §7, No.~1, Prop. 1)，于是 \(p_{\mathbf{K}} = (\mu_{\mathbf{K}})^{\bullet}\) 。因而 Prop. 2 b) 的条件 1) 也得到满足，\(\mu\) 的存在性随之立得。

证毕。

我们称 \(\mu\) 为由 \(\lambda\) 定义的 T 上的测度。特别地，对 T 的每一点 \(x\) 可以定义测度 \(\varepsilon_x\) ；它由 \((\varepsilon_x)^{\bullet}(f) = f(x)\) （对 \(f \in \mathcal{F}_{+}(\mathrm{T})\)）刻画。

注记. --- 1) 当 T 局部紧时，\(\mu\) 是 \(\lambda\) 在 L 到 T 的单射下的像。我们将在 §2, No.~3 的例中（在处理过像测度之后）看到，这一解释对任意空间仍然有效。

\begin{enumerate}
\def\labelenumi{\arabic{enumi})}
\setcounter{enumi}{1}
\tightlist
\item
  我们还将看到，例 2 中定义的测度是 T 上具有紧支集的正测度 (No.~6, 注 2)。
\end{enumerate}

除明确说明相反情形外，今后我们只考虑正测度。在本节余下部分，T 表示一个拓扑空间，\(\mu\) 表示 T 上的一个正测度。

以下各小节中的许多结果可以推广到正前测度。这一推广留给读者。

